% =============================================================================
\section{Méthodes explorées}
\label{sec:autres_algorithmes}
% =============================================================================

Au-delà des trois familles retenues (section~\ref{sec:methods}), nous avons
implémenté et évalué un ensemble étendu de variantes et d'algorithmes
complémentaires.  Cette exploration systématique a permis de valider le choix
final par élimination et de quantifier l'apport de chaque mécanisme.  Nous
présentons ici ces méthodes de façon synthétique.

\subsection{Algorithme génétique réel (GA) — ligne de base}
\label{sec:ga}

L'algorithme génétique réel~\cite{deb2001} sert de \textit{baseline} classique.
Il utilise un croisement SBX (\textit{Simulated Binary Crossover},
$\eta_c = 20$, $p_c = 0{,}9$), une mutation polynomiale ($\eta_m = 20$,
$p_m = 1/d$), une sélection par tournoi binaire et un élitisme du meilleur
individu.  La population est fixée à $N_P = 100$.

Sur nos instances, le GA atteint des scores comparables à CMA-ES sur
\texttt{prob2} ($d=8$) mais se dégrade rapidement sur les problèmes corrélés
(\texttt{prob3} : $42{,}1$ vs $28{,}5$ pour CMA-ES) en raison du croisement
composante par composante qui ne capture pas les corrélations entre variables.
Ce résultat justifie le recours à des méthodes adaptant explicitement la
structure de corrélation.

\subsection{Variantes de l'évolution différentielle}

\paragraph{SHADE (\textit{Success-History based Adaptive DE})~\cite{tanabe2013}.}
Au lieu d'adapter $F$ et $CR$ individuellement comme jDE, SHADE maintient une
mémoire circulaire de taille $H = 10$ stockant les paramètres ayant conduit à
des améliorations.  Les nouveaux $F$ sont tirés d'une distribution de Cauchy et
les $CR$ d'une gaussienne, centrés sur la moyenne de Lehmer des succès
historiques.  Sur \texttt{prob3}, SHADE atteint $28{,}45 \pm 0{,}31$, égalant
CMA-ES.

\paragraph{L-SHADE (\textit{Linear Population Size Reduction})~\cite{tanabe2014}.}
Extension de SHADE avec une réduction linéaire de la population de $N_0 = 18d$
à $N_{\min} = 4$ individus, combinée à une archive externe des solutions
remplacées.  La mutation utilise \texttt{current-to-pbest/1} avec $p = 0{,}11$.

\subsection{Variantes PSO}

\paragraph{CLPSO (\textit{Comprehensive Learning PSO}).}
Chaque dimension d'une particule apprend d'un exemplaire potentiellement
différent, sélectionné par tournoi avec une probabilité $P_c(i)$ croissante
selon le rang.  L'essaim est de 40~particules avec un intervalle de
rafraîchissement de 7~générations.

\paragraph{FPSO (\textit{Fractional-Order PSO}).}
Variante intégrant les dérivées de Grünwald-Letnikov d'ordre fractionnaire
$\beta \in [0{,}4, \; 0{,}9]$ (décroissant) sur une fenêtre mémoire de $M = 7$
vitesses passées.

\subsection{Algorithmes d'intelligence en essaim}

\paragraph{GWO (\textit{Grey Wolf Optimizer}).}
Hiérarchie de meute (alpha, bêta, delta) guidant 30~loups, avec un paramètre
d'exploration $a$ décroissant linéairement de 2 à 0.

\paragraph{FA (\textit{Firefly Algorithm}).}
Algorithme des lucioles (30~individus) avec une attractivité $\beta =
\beta_0 \exp(-\gamma r^2)$ décroissante avec la distance.

\subsection{Hybridation avec la planification discrète}

\paragraph{AStarSeeds.}
L'algorithme A*~\cite{hart1968} fournit un chemin discret sans collision sur
une grille de navigation.  Les points de contrôle de la courbe de Bézier sont
initialisés le long de ce chemin (\textit{seeding}), puis l'optimiseur continu
(CMA-ES) affine la trajectoire.

\paragraph{AStarRepair.}
En cours d'optimisation, les individus traversant un obstacle sont réparés
localement par un appel à A*.  Cette approche atteint le meilleur score sur
\texttt{prob3} ($28{,}24 \pm 1{,}52$).

\subsection{Gestion avancée des contraintes}

\paragraph{ALConstraint (\textit{Augmented Lagrangian}).}
Transformation du problème pénalisé en sous-problèmes sans contrainte avec
mise à jour dynamique des multiplicateurs de Lagrange.

\paragraph{StochRank (\textit{Stochastic Ranking}).}
Lors du tri de la population, la \textit{fitness} et le niveau de violation sont
comparés avec une probabilité $P_f = 0{,}45$.  Cette relaxation permet à des
individus légèrement infaisables mais prometteurs de survivre, favorisant le
contournement fluide des obstacles.

\subsection{Autres variantes CMA-ES}

\paragraph{LMCMA (\textit{Limited Memory CMA-ES}).}
Approximation de rang faible de la covariance réduisant la complexité à
$\mathcal{O}(d \log d)$, conçue pour $d > 1000$.  Sur nos instances ($d \leq 32$),
LMCMA s'est montré instable (\texttt{prob3} : $69{,}4 \pm 23{,}5$).

\paragraph{RFSurrogate (\textit{Random Forest Surrogate}).}
Pré-évaluation des candidats par un modèle de forêts aléatoires.  Seuls les
individus jugés prometteurs subissent l'évaluation exacte.  Bon sur
\texttt{prob3} ($28{,}35 \pm 0{,}20$) mais ajoutant de la variance.

\paragraph{Islands, GridBIPOP, SepWarmup, MultiSegment.}
Le modèle insulaire (migration inter-sous-populations), le BIPOP couplé à un
archive MAP-Elites, le démarrage en covariance diagonale puis pleine, et la
spline multi-segments ont été testés sans apporter de gain significatif sur
l'ensemble des quatre instances.

\subsection{Bilan de l'exploration}

Les tableaux~\ref{tab:results_prob123} et~\ref{tab:results_prob4}
(section~\ref{sec:experiments}) montrent que
la grande majorité de ces variantes offrent des performances comparables ou
inférieures aux trois méthodes retenues.  Les apports les plus notables sont :
\begin{itemize}[nosep]
  \item SHADE, qui égale CMA-ES sur \texttt{prob3} grâce à l'adaptation
        historique de $F$ et $CR$ ;
  \item AStarRepair, qui fournit le meilleur score absolu sur \texttt{prob3}
        via l'hybridation planification discrète + optimisation continue ;
  \item StochRank et BIPOP, seules variantes CMA-ES à descendre sous 5\,000
        sur \texttt{prob4}.
\end{itemize}
Ces observations confirment la pertinence de concentrer les efforts futurs sur
CMA-ES (avec \textit{restarts} et bornes étendues), PSO et DE/SHADE.
